% Options for packages loaded elsewhere
% Options for packages loaded elsewhere
\PassOptionsToPackage{unicode}{hyperref}
\PassOptionsToPackage{hyphens}{url}
\PassOptionsToPackage{dvipsnames,svgnames,x11names}{xcolor}
%
\documentclass[
  letterpaper,
  DIV=11,
  numbers=noendperiod]{scrartcl}
\usepackage{xcolor}
\usepackage{amsmath,amssymb}
\setcounter{secnumdepth}{-\maxdimen} % remove section numbering
\usepackage{iftex}
\ifPDFTeX
  \usepackage[T1]{fontenc}
  \usepackage[utf8]{inputenc}
  \usepackage{textcomp} % provide euro and other symbols
\else % if luatex or xetex
  \usepackage{unicode-math} % this also loads fontspec
  \defaultfontfeatures{Scale=MatchLowercase}
  \defaultfontfeatures[\rmfamily]{Ligatures=TeX,Scale=1}
\fi
\usepackage{lmodern}
\ifPDFTeX\else
  % xetex/luatex font selection
\fi
% Use upquote if available, for straight quotes in verbatim environments
\IfFileExists{upquote.sty}{\usepackage{upquote}}{}
\IfFileExists{microtype.sty}{% use microtype if available
  \usepackage[]{microtype}
  \UseMicrotypeSet[protrusion]{basicmath} % disable protrusion for tt fonts
}{}
\makeatletter
\@ifundefined{KOMAClassName}{% if non-KOMA class
  \IfFileExists{parskip.sty}{%
    \usepackage{parskip}
  }{% else
    \setlength{\parindent}{0pt}
    \setlength{\parskip}{6pt plus 2pt minus 1pt}}
}{% if KOMA class
  \KOMAoptions{parskip=half}}
\makeatother
% Make \paragraph and \subparagraph free-standing
\makeatletter
\ifx\paragraph\undefined\else
  \let\oldparagraph\paragraph
  \renewcommand{\paragraph}{
    \@ifstar
      \xxxParagraphStar
      \xxxParagraphNoStar
  }
  \newcommand{\xxxParagraphStar}[1]{\oldparagraph*{#1}\mbox{}}
  \newcommand{\xxxParagraphNoStar}[1]{\oldparagraph{#1}\mbox{}}
\fi
\ifx\subparagraph\undefined\else
  \let\oldsubparagraph\subparagraph
  \renewcommand{\subparagraph}{
    \@ifstar
      \xxxSubParagraphStar
      \xxxSubParagraphNoStar
  }
  \newcommand{\xxxSubParagraphStar}[1]{\oldsubparagraph*{#1}\mbox{}}
  \newcommand{\xxxSubParagraphNoStar}[1]{\oldsubparagraph{#1}\mbox{}}
\fi
\makeatother


\usepackage{longtable,booktabs,array}
\usepackage{calc} % for calculating minipage widths
% Correct order of tables after \paragraph or \subparagraph
\usepackage{etoolbox}
\makeatletter
\patchcmd\longtable{\par}{\if@noskipsec\mbox{}\fi\par}{}{}
\makeatother
% Allow footnotes in longtable head/foot
\IfFileExists{footnotehyper.sty}{\usepackage{footnotehyper}}{\usepackage{footnote}}
\makesavenoteenv{longtable}
\usepackage{graphicx}
\makeatletter
\newsavebox\pandoc@box
\newcommand*\pandocbounded[1]{% scales image to fit in text height/width
  \sbox\pandoc@box{#1}%
  \Gscale@div\@tempa{\textheight}{\dimexpr\ht\pandoc@box+\dp\pandoc@box\relax}%
  \Gscale@div\@tempb{\linewidth}{\wd\pandoc@box}%
  \ifdim\@tempb\p@<\@tempa\p@\let\@tempa\@tempb\fi% select the smaller of both
  \ifdim\@tempa\p@<\p@\scalebox{\@tempa}{\usebox\pandoc@box}%
  \else\usebox{\pandoc@box}%
  \fi%
}
% Set default figure placement to htbp
\def\fps@figure{htbp}
\makeatother





\setlength{\emergencystretch}{3em} % prevent overfull lines

\providecommand{\tightlist}{%
  \setlength{\itemsep}{0pt}\setlength{\parskip}{0pt}}



 


\KOMAoption{captions}{tableheading}
\makeatletter
\@ifpackageloaded{caption}{}{\usepackage{caption}}
\AtBeginDocument{%
\ifdefined\contentsname
  \renewcommand*\contentsname{Table of contents}
\else
  \newcommand\contentsname{Table of contents}
\fi
\ifdefined\listfigurename
  \renewcommand*\listfigurename{List of Figures}
\else
  \newcommand\listfigurename{List of Figures}
\fi
\ifdefined\listtablename
  \renewcommand*\listtablename{List of Tables}
\else
  \newcommand\listtablename{List of Tables}
\fi
\ifdefined\figurename
  \renewcommand*\figurename{Figure}
\else
  \newcommand\figurename{Figure}
\fi
\ifdefined\tablename
  \renewcommand*\tablename{Table}
\else
  \newcommand\tablename{Table}
\fi
}
\@ifpackageloaded{float}{}{\usepackage{float}}
\floatstyle{ruled}
\@ifundefined{c@chapter}{\newfloat{codelisting}{h}{lop}}{\newfloat{codelisting}{h}{lop}[chapter]}
\floatname{codelisting}{Listing}
\newcommand*\listoflistings{\listof{codelisting}{List of Listings}}
\makeatother
\makeatletter
\makeatother
\makeatletter
\@ifpackageloaded{caption}{}{\usepackage{caption}}
\@ifpackageloaded{subcaption}{}{\usepackage{subcaption}}
\makeatother
\usepackage{bookmark}
\IfFileExists{xurl.sty}{\usepackage{xurl}}{} % add URL line breaks if available
\urlstyle{same}
\hypersetup{
  pdftitle={STA5092Z 2026 Assignment 1},
  pdfauthor={Pamela Afful},
  pdfkeywords={pied babbler *comeback and edit properly},
  colorlinks=true,
  linkcolor={blue},
  filecolor={Maroon},
  citecolor={Blue},
  urlcolor={Blue},
  pdfcreator={LaTeX via pandoc}}


\title{STA5092Z 2026 Assignment 1}
\author{Pamela Afful}
\date{}
\begin{document}
\maketitle
\begin{abstract}
I will doing some EDA on very messy data, yay!
\end{abstract}


\section{Part 1- Data Wrangling}\label{part-1--data-wrangling}

\subsubsection{Load Data}\label{load-data}

2) Load data and transform into tibble objects

\begin{verbatim}
# A tibble: 596 x 1
   Date     
   <chr>    
 1 3-Oct-05 
 2 3-Oct-05 
 3 3-Oct-05 
 4 4-Oct-05 
 5 4-Oct-05 
 6 12-Oct-05
 7 12-Oct-05
 8 12-Oct-05
 9 15-Oct-05
10 15-Oct-05
# i 586 more rows
\end{verbatim}

\begin{verbatim}
# A tibble: 427 x 1
   Date     
   <chr>    
 1 03-Oct-05
 2 03-Oct-05
 3 04-Oct-05
 4 04-Oct-05
 5 12-Oct-05
 6 12-Oct-05
 7 12-Oct-05
 8 15-Oct-05
 9 15-Oct-05
10 15-Oct-05
# i 417 more rows
\end{verbatim}

\subsubsection{Merge Tibble Dataframes}\label{merge-tibble-dataframes}

3)

\subsubsection{DataFrame Column Data
Types}\label{dataframe-column-data-types}

\paragraph{Key Findings:}\label{key-findings}

\textbf{fledg\_survival} dataframe

\begin{itemize}
\item
  Data types in columns are mostly split between character and
  int/numeric data types
\item
  \textbf{Date} column in fledg\_survival is stored as a character- will
  need to be converted. The data format is of the format dd-mm-yy
\item
  \textbf{SurvMassAtInd} column contains null values
\end{itemize}

\textbf{nest\_all} dataframe

\begin{itemize}
\item
  In addition to numeric and int/numeric data types, columns also
  include logical/boolean data types
\item
  \textbf{Date} is also a character data type, similar to the format in
  the \textbf{fledg\_survival} dataframe
\end{itemize}

Merging steps:

\begin{enumerate}
\def\labelenumi{\arabic{enumi}.}
\item
  First, convert date columns in both data frames into datetime columns
  using the lubridate ymd function
\item
  Find common column names in both the \textbf{nest\_all} and
  \textbf{fledg\_survival} data sets to find columns to join on
\item
  Check and ensure that the datatypes of these common columns are the
  same in both data sets- using the janitor package- and change any
  datatype mismatches to the same data type
\item
  Both GrpSizeAD and GrpSizeTotal columns are stored as characters in
  the nest\_all data frame, as they both contain \#NA values
\item
  These will replaced with na values instead, once they are coerced into
  integer values
\end{enumerate}

\section{Part 2- Data Summaries and
Visualizations}\label{part-2--data-summaries-and-visualizations}

1) \textbf{piedbabbler} samplesize

\begin{verbatim}
[1] 862
\end{verbatim}

2)

\paragraph{Null Value Check}\label{null-value-check}

\begin{verbatim}
# A tibble: 1 x 8
  GrpSizeAD GrpSizeTotal BroodSize  Mass TMaxMeas No.ChicksFledge Rain90 Drought
      <int>        <int>     <int> <int>    <int>           <int>  <int>   <int>
1         1            1         0     8      303             267    159      26
\end{verbatim}

35\% of values i

\paragraph{Outlier Analysis}\label{outlier-analysis}

Outlier Definition: An outlier is an individual point of data that is
distant from other points in the data set. It is an anomaly in the
dataset that may be caused by a range of errors in capturing, processing
or manipulating data. * comeback and reference

\pandocbounded{\includegraphics[keepaspectratio]{EDA_Assignment1_files/figure-pdf/out_lier-1.pdf}}

From their respective boxplots, the following columns seem to contain
outliers, relative to their own distributions:

\begin{itemize}
\item
  TMaxMeas
\item
  Mass
\item
  BroodSize
\item
  GrpSizeTotal
\end{itemize}

3)

\paragraph{Mass and Date per Nest: Longitudinal
Dataframe}\label{mass-and-date-per-nest-longitudinal-dataframe}

\begin{verbatim}
# A tibble: 257 x 3
# Groups:   NestID [257]
   NestID Date       total_mass
   <chr>  <date>          <dbl>
 1 AQAN1C 2005-10-22      118. 
 2 AQAN1D 2006-10-14       74.7
 3 AQAN2C 2005-12-19       46  
 4 BABN1D 2006-11-08       65.2
 5 BABN3C 2005-11-28      127. 
 6 BABN4C 2006-02-11      191. 
 7 BALN2C 2005-11-09       79.5
 8 BALN3C 2006-02-18       98  
 9 BMDN1J 2011-10-29       86.9
10 BMDN1K 2012-11-21       98  
# i 247 more rows
\end{verbatim}

4)

TMaxMeas has 35\% of its values as nulls. As such, the aggregate mean
Mass value per bucketed TMaxMeas will be use to adjust for this

\pandocbounded{\includegraphics[keepaspectratio]{EDA_Assignment1_files/figure-pdf/crit_tem-1.pdf}}

5)

Mass over time per nest per season

Since we have way too distinct NestIDs, the following grouping approach
will be used:\\

\begin{enumerate}
\def\labelenumi{\arabic{enumi}.}
\item
  First get the size of each nest, defined by the count of birds in each
  nest- since this value, changes over time, the starting nest size for
  each group will be used
\item
  Use this count as the grouping variable to compare mass over time
\item
  Show the individual nestID groups for the ``best'' performing birds
\item
  Also, since there are not enough observations for each nest ID, per
  year, the year variable will be normalized relatve to the start
  date/year of the nest
\end{enumerate}

\begin{verbatim}
# A tibble: 14 x 3
   Season num_index num_index_break
   <chr>      <int> <fct>          
 1 C              1 1              
 2 D              2 1              
 3 E              3 1              
 4 F              4 1              
 5 G              5 1              
 6 H              6 2              
 7 J              7 2              
 8 K              8 2              
 9 L              9 2              
10 M             10 3              
11 N             11 3              
12 P             12 3              
13 Q             13 3              
14 R             14 3              
\end{verbatim}

\begin{verbatim}
# A tibble: 862 x 87
   Season num_index num_index_break Date       BirdID Group AgeMeasD Sex  
   <chr>      <int> <fct>           <date>     <chr>  <chr>    <int> <chr>
 1 C              1 1               2005-10-03 MLYG   HAR         11 F    
 2 C              1 1               2005-10-03 PMPT   HAR         11 F    
 3 C              1 1               2005-10-04 WRMG   XHO         13 F    
 4 C              1 1               2005-10-04 RMPR   XHO         13 M    
 5 C              1 1               2005-10-12 PPRM   TPT         11 F    
 6 C              1 1               2005-10-12 TRMP   TPT         11 F    
 7 C              1 1               2005-10-12 LMBW   TPT         11 M    
 8 C              1 1               2005-10-15 BGPM   OSD         11 F    
 9 C              1 1               2005-10-15 MBBR   OSD         11 M    
10 C              1 1               2005-10-15 PTMR   OSD         11 M    
# i 852 more rows
# i 79 more variables: NestID <chr>, GrpSizeAD <int>, GrpSizeTotal <int>,
#   PairTenure <int>, BroodSize <int>, NoMales <dbl>, BroodSexRatio <dbl>,
#   Mass <dbl>, SurvInd <int>, SurvPeriod <int>, SurvMassAtInd <dbl>,
#   Surv1Yr <int>, Surv1Mass <chr>, MeanTmaxInc <dbl>, MeanTmaxBrood <dbl>,
#   MeanTmaxInd <dbl>, MeanTminInc <dbl>, MeanTminBrood <dbl>,
#   MeanTmin90 <dbl>, MeanTvarInc <dbl>, MeanTvarBrood <dbl>, ...
\end{verbatim}

\pandocbounded{\includegraphics[keepaspectratio]{EDA_Assignment1_files/figure-pdf/plt_c_g-1.pdf}}

\pandocbounded{\includegraphics[keepaspectratio]{EDA_Assignment1_files/figure-pdf/plt_h_j-1.pdf}}

\pandocbounded{\includegraphics[keepaspectratio]{EDA_Assignment1_files/figure-pdf/plt_m_r-1.pdf}}

Although the time ranges are inconsistent and disparate across seasons,
nest sizes with 3 chicks, show higher masses over time. This will be
used as the criterion to filter and observe individual nest ids that
fall in this category.

\pandocbounded{\includegraphics[keepaspectratio]{EDA_Assignment1_files/figure-pdf/mess_ev-1.pdf}}

6)

\textbf{Yearly Average Mass, Faceted by Season}

\pandocbounded{\includegraphics[keepaspectratio]{EDA_Assignment1_files/figure-pdf/chick_mass_drought-1.pdf}}

\textbf{Average Bird Size Over Time- Faceted By Season}

\pandocbounded{\includegraphics[keepaspectratio]{EDA_Assignment1_files/figure-pdf/chick_size_drought-1.pdf}}

\pandocbounded{\includegraphics[keepaspectratio]{EDA_Assignment1_files/figure-pdf/grp_size_drought-1.pdf}}

\pandocbounded{\includegraphics[keepaspectratio]{EDA_Assignment1_files/figure-pdf/chick_fledge_drought-1.pdf}}

\pandocbounded{\includegraphics[keepaspectratio]{EDA_Assignment1_files/figure-pdf/rain_drought-1.pdf}}

6) Rate of mass change per day per nesting

Steps for finding these suspicious entries

\begin{enumerate}
\def\labelenumi{\arabic{enumi}.}
\tightlist
\item
  First group Mass by Group and Date to get Mass per Nest over time
\item
  Get lagged values of Date and Mass
\item
  Calculate the percentage Mass change for each row per nest group
\item
  Calculate the date difference, in days, between Date values and lagged
  date values
\item
  Filter for date differences equal to 1 then check which growth rates
  are suspicious (\textgreater50\%)
\end{enumerate}

\begin{verbatim}
# A tibble: 1 x 1
# Groups:   Group [1]
  Group
  <chr>
1 <NA> 
\end{verbatim}

All these above values are NAs*, so I guess there was no data for these
values(?)

8)

\pandocbounded{\includegraphics[keepaspectratio]{EDA_Assignment1_files/figure-pdf/mass_fleg_ind-1.pdf}}

\pandocbounded{\includegraphics[keepaspectratio]{EDA_Assignment1_files/mediabag/000020.png}}

\pandocbounded{\includegraphics[keepaspectratio]{EDA_Assignment1_files/figure-pdf/mass_diff-1.pdf}}

9)

First check how the number of adult helpers are distributed across the
data

\pandocbounded{\includegraphics[keepaspectratio]{EDA_Assignment1_files/figure-pdf/ad_help-1.pdf}}

In this case, most nests had adult heplers of 3. Thus, to normalize the
data, average fledgling count per adult helper count we'll be use to
compare.

9)

\pandocbounded{\includegraphics[keepaspectratio]{EDA_Assignment1_files/figure-pdf/help_cnt-1.pdf}}

The average number fledged chicks sees an increase from 3 to 6, however
it sees its most marked sustained decrease from 6 to 8 ( although this
rebounds after 8)

\pandocbounded{\includegraphics[keepaspectratio]{EDA_Assignment1_files/figure-pdf/cum_per-1.pdf}}

10)

Comeback*

Nestiling Mass Per Drought Per Nest Size

Define big nest size as GrpSizeTotal \textgreater6

Small nest size GrpSizeTotal \textless6

\pandocbounded{\includegraphics[keepaspectratio]{EDA_Assignment1_files/figure-pdf/drough_comp-1.pdf}}

Smaller nests seem to survive longer over periods of drought*

11)

Seasonal Heatmap

We need to bucket the ambient values

Nesting mortality=number of dead birds/total number of birds

Come back and explain

\pandocbounded{
\includegraphics[keepaspectratio]{EDA_Assignment1_files/figure-pdf/heat_map-1.pdf}}

12)

Ridge Plot

\pandocbounded{\includegraphics[keepaspectratio]{EDA_Assignment1_files/figure-pdf/ridge-1.pdf}}

\textbf{Season Q} has the most bimodal distribution

13) Comeback and think if a nice plot- Sankey Plot?

\pandocbounded{\includegraphics[keepaspectratio]{EDA_Assignment1_files/figure-pdf/unnamed-chunk-3-1.pdf}}




\end{document}
